\chapter{Methodology}

This chapter describes the data sources, variables, causal framework, and statistical methods used to estimate the causal effect of ITN use on malaria infection.

\section{Data Sources}

This study uses data from the Kenya Malaria Indicator Surveys (KMIS) conducted in 2015 and 2020 \citep{kmis2015report, kmis2020report}. The KMIS is a nationally representative household survey that collects information on malaria prevention, diagnosis, and treatment. Unlike the Demographic and Health Surveys (DHS) which do not include malaria parasite testing, the KMIS specifically focuses on malaria and includes biomarker data for malaria infection.

\subsection{KMIS 2015 and 2020}

The 2015 KMIS was conducted between July and December 2015, selecting 245 enumeration areas (clusters) and interviewing 6,481 households, including 3,925 children under five. The 2020 KMIS was conducted between February and December 2020, selecting 298 clusters and interviewing 7,952 households, including 4,298 children under five.

\subsection{Analytical Sample Construction}

Starting from the full KMIS samples, we applied the following steps: (1) filter to children under five years of age (age < 60 months), (2) merge household-level and environmental data, (3) exclude observations with missing outcome or treatment variables, and (4) exclude observations with missing covariate data. The final analytical sample consists of 6,771 children across 528 unique clusters in Kenya. Malaria prevalence in the combined sample is 4.8\% (326 positive cases), and the ITN use rate is 54.1\%.

\subsection{Geographic and Environmental Data}

Environmental variables were extracted from the DHS Geocov datasets \citep{dhsgeocov}, which link survey clusters to high-resolution environmental data: rainfall (annual precipitation in mm), temperature (mean annual temperature in °C), Enhanced Vegetation Index (EVI), aridity index, wet days per year, population count, and elevation (2020 only).

\section{Variables and Measurement}

\subsection{Outcome Variable: Malaria Infection}

The outcome variable is malaria infection status among children under five, recorded in variable \texttt{hml32} as the result of a rapid diagnostic test (RDT) or microscopy. The variable is coded as 1 if the child tests positive for malaria and 0 if negative. 

\subsection{Treatment Variable: ITN Use}

The treatment variable is ITN use, measured by whether the child slept under an ITN the night before the survey, recorded in variable \texttt{hml12}. The coding is 1 if the child slept under an ITN and 0 otherwise. This measure captures actual protective exposure rather than simple net ownership.

\subsection{Covariates}

Covariates include demographic factors (age in months, sex), socioeconomic factors (wealth quintile, maternal education, urban/rural residence), and environmental factors (rainfall, temperature, EVI, aridity, wet days, population, elevation).

\subsection{Missing Data Handling}

Observations with missing outcome or treatment variables are excluded. Covariate missingness is minimal (less than 5\% for all covariates), and complete case analysis is employed following \citet{little2019statistical}.

\section{Causal Framework and Identification Strategy}

\subsection{The Potential Outcomes Framework}

Following the Rubin Causal Model \citep{rubin1974estimating}, we define potential outcomes for each child $i$. Let $Y_i(1)$ be the malaria infection status if the child uses an ITN, and $Y_i(0)$ be the status if the child does not use an ITN. The individual treatment effect is $\tau_i = Y_i(1) - Y_i(0)$. The fundamental problem of causal inference is that we only observe one potential outcome for each child: $Y_i^{\text{obs}} = T_i Y_i(1) + (1 - T_i) Y_i(0)$, where $T_i$ is the observed treatment indicator.

\subsection{Parameter of Interest}

The primary parameter of interest is the Average Treatment Effect (ATE): $\text{ATE} = \mathbb{E}[Y(1) - Y(0)]$. We estimate two types of effects: cluster-specific (conditional) effects using GLMM-based methods, and population-average (marginal) effects using DML.

\subsection{Identification Assumptions}

To identify the ATE from observational data, we require: (1) conditional ignorability: $\{Y(1), Y(0)\} \perp T \mid \mathbf{X}$; (2) positivity: $0 < P(T = 1 \mid \mathbf{X}) < 1$ for all $\mathbf{X}$ \citep{westreich2010invited}; (3) consistency; and (4) no interference (SUTVA). The balancing property of the propensity score \citep{rosenbaum1983central} ensures that $T \perp \mathbf{X} \mid e(\mathbf{X})$.

\subsection{Directed Acyclic Graph}

Figure~\ref{fig:dag} presents the Directed Acyclic Graph (DAG) encoding our causal assumptions. The DAG shows the direct causal path (ITN use $\rightarrow$ Malaria), confounders (age, sex, wealth, maternal education, residence, and environmental factors), and no unblocked backdoor paths after conditioning on the identified confounders.

\begin{figure}[htbp]
	\centering
	\includegraphics[width=0.9\textwidth]{Figures/dag_causal_diagram_ggraph.png}
	\caption{Directed Acyclic Graph (DAG) for ITN Use and Malaria}
	\label{fig:dag}
\end{figure}

\section{Statistical Methods}

\subsection{Baseline Regression and GLMM}

For child $i$ in cluster $c$, the probability of malaria infection $p_{ic} = P(Y_{ic} = 1 \mid \mathbf{X}_{ic})$ is modeled using survey-weighted logistic regression:

\[
\log\left(\frac{p_{ic}}{1-p_{ic}}\right) = \beta_0 + \beta_1 \text{ITN}_{ic} + \boldsymbol{\gamma}^\top \mathbf{X}_{ic}
\]

We estimate three specifications: unadjusted (no covariates), demographic + socioeconomic (age, sex, wealth, residence), and fully adjusted (all covariates from the DAG).

To account for clustering, we employ GLMM with a random intercept $u_c$ for each cluster:

\[
\log\left(\frac{p_{ic}}{1-p_{ic}}\right) = \mathbf{X}_{ic}^\top \boldsymbol{\beta} + u_c, \quad u_c \sim N(0, \sigma_u^2)
\]

The intra-cluster correlation (ICC) measures the proportion of total variance explained by between-cluster differences:

\[
\text{ICC} = \frac{\sigma_u^2}{\sigma_u^2 + \pi^2/3}
\]

Following \citet{killip2004intracluster}, ICC values above 0.05 justify the use of random effects models. Our preliminary analysis reveals ICC = 41-47\%.

\subsection{Propensity Score Methods with GLMM}

The propensity score \citep{rosenbaum1983central} is $e(\mathbf{X}) = P(T = 1 \mid \mathbf{X})$. To account for clustering, we estimate propensity scores using a GLMM:

\[
\log\left(\frac{e(\mathbf{X}_{ic})}{1 - e(\mathbf{X}_{ic})}\right) = \alpha_0 + \boldsymbol{\alpha}^\top \mathbf{X}_{ic} + v_c, \quad v_c \sim N(0, \sigma_v^2)
\]

We implement three methods: (1) Propensity Score Matching (PSM) using nearest neighbor matching with a caliper of 0.2 standard deviations \citep{austin2011introduction, austin2010optimal}; (2) Inverse Probability of Treatment Weighting (IPTW) using stabilized weights \citep{cole2008constructing}; and (3) Doubly Robust estimation \citep{kang2007average}.

\subsection{Double Machine Learning Framework}

We use DML2 (Definition 3.2 in \citealp{chernozhukov2018double}) with the partially linear model \citep{nie2021quasi}:

\[
Y = \theta D + g(X) + \varepsilon, \quad D = m(X) + \eta
\]

where $Y$ is malaria infection, $D$ is ITN use, $\theta$ is the ATE, and $g(X)$ and $m(X)$ are nuisance functions estimated using machine learning. We employ two algorithms: Random Forest \citep{breiman2001random} and XGBoost \citep{chen2016xgboost}, with 5-fold and 10-fold cross-fitting.

\subsection{Heterogeneity Analysis}

We estimate subgroup-specific treatment effects using DML for wealth quintiles, residence type (urban/rural), age groups (0-11, 12-23, 24-35, 36-59 months), and rainfall levels (low, medium, high based on tertiles).